\chapter{Conexión de la Aplicación Android}

% Texto de sección mejorado basado en el comportamiento del repositorio MICSBOL/ESP32_BT_Controller.
% Notas:
% - Se usa \verb|...| para nombres de archivos/macros y así evitar problemas con guiones bajos en babel+hyperref.
% - Las banderas de depuración reflejan la configuración en debug_config.h del repositorio.

\section{Instalar la Aplicación}
Instala la aplicación Android de RC Controller en tu teléfono usando el método de distribución proporcionado
por tu proyecto (APK, enlace de tienda o despliegue interno).

\section{Activar Bluetooth}
Activa Bluetooth en tu teléfono antes de abrir la aplicación. En versiones recientes de Android, la app puede solicitar
permisos adicionales (por ejemplo, ``Dispositivos cercanos'' o permisos relacionados con Ubicación) para poder escanear
dispositivos Bluetooth. Concede los permisos solicitados para que la lista de dispositivos se muestre correctamente.

\section{Conectarse al ESP32}
Enciende el ESP32 y abre la aplicación RC Controller. En la lista de dispositivos Bluetooth disponibles, selecciona
el controlador ESP32.

Una vez establecida la conexión, la interfaz de control se activa y la aplicación empieza a enviar datos de control
al ESP32 (joysticks, perillas, interruptores y eventos de botones).

\section{Estado de Conexión}
La aplicación muestra el estado de conexión para confirmar que el enlace serie Bluetooth está activo.

Del lado del firmware, un cliente conectado se detecta con \verb|SerialBT.hasClient()| (utilizado por el módulo
de telemetría). Si el teléfono se desconecta, el firmware vuelve al estado de ``espera'' y la aplicación debe reconectar
antes de que el control y la telemetría se reanuden.

\section{Modo Debug (Ver la Comunicación Bluetooth)}
El firmware incluye un sistema de depuración opcional que imprime datos del controlador ya decodificados en el Monitor
Serie USB. Esta es la forma más rápida de confirmar que:
\begin{itemize}
    \item el teléfono está conectado al dispositivo ESP32 correcto,
    \item los paquetes se están recibiendo y decodificando correctamente, y
    \item los valores de sticks/perillas/interruptores cambian cuando mueves los controles en la app.
\end{itemize}

\subsection{Activar el Modo Debug}
Abre el archivo:

\begin{verbatim}
    debug_config.h
\end{verbatim}

Activa la depuración estableciendo:

\begin{verbatim}
    #define DEBUG_ENABLED 1
\end{verbatim}

Desactiva la depuración estableciendo:

\begin{verbatim}
    #define DEBUG_ENABLED 0
\end{verbatim}

Cuando \verb|DEBUG_ENABLED| está activo, el repositorio también habilita categorías individuales de depuración:

\begin{verbatim}
    #define DBG_STICKS   1
    #define DBG_KNOBS    1
    #define DBG_SWITCHES 1
    #define DBG_EVENTS   1
\end{verbatim}

Después de cambiar cualquier ajuste de depuración, compila y carga el firmware nuevamente en el ESP32.

\subsection{Abrir el Monitor Serie}

Abre el Monitor Serie del IDE de Arduino:

\begin{verbatim}
    Tools -> Serial Monitor
\end{verbatim}

Configura la velocidad (baud rate) en:

\begin{verbatim}
    115200
\end{verbatim}

Reinicia el ESP32 para ver los mensajes de arranque.

\subsection{Salida de Depuración}

Cuando el modo debug está habilitado, el ESP32 imprimirá mensajes de diagnóstico en el Monitor Serie.

Mensajes típicos pueden incluir:

\begin{verbatim}
    [DEBUG] Bluetooth initialized
    [DEBUG] Waiting for client connection
    [DEBUG] Client connected
\end{verbatim}

Cuando se reciben comandos desde la aplicación Android, el ESP32 puede mostrar mensajes similares a:

\begin{verbatim}
    [DEBUG] Received command: BTN_A
    [DEBUG] Joystick X: 120
    [DEBUG] Joystick Y: 85
\end{verbatim}

Estos mensajes confirman que la comunicación Bluetooth entre el teléfono y el ESP32 está funcionando correctamente.

\subsection{Solución de Problemas Usando el Modo Debug}

Si la aplicación Android se conecta pero el robot no responde:

\begin{itemize}
    \item Verifica que aparezcan mensajes de depuración al presionar botones en la aplicación.
    \item Si no aparece ningún mensaje, es posible que el ESP32 no esté recibiendo datos Bluetooth.
    \item Comprueba que en la aplicación se haya seleccionado el dispositivo ESP32 correcto.
    \item Verifica que el firmware se haya cargado correctamente.
\end{itemize}

Si aparecen mensajes de depuración pero el hardware no responde, el problema probablemente está relacionado con el cableado o con la configuración de pines.

\section{Usar el Modo Debug para Verificar la Comunicación}

El firmware incluye un sistema de depuración integrado que permite ver los datos que se reciben desde la aplicación Android.

Esto es extremadamente útil para comprobar que la conexión Bluetooth funciona correctamente y que el ESP32 está recibiendo comandos de control.

Todos los mensajes de depuración se muestran en el Monitor Serie del IDE de Arduino.

\subsection{Paso 1: Activar el Modo Debug}

Abre el archivo:

\begin{verbatim}
    debug_config.h
\end{verbatim}

Ubica la siguiente línea:

\begin{verbatim}
    #define DEBUG_ENABLED 1
\end{verbatim}

Ajusta el valor a:

\begin{itemize}
    \item \texttt{1} para habilitar la salida de depuración
    \item \texttt{0} para deshabilitar la salida de depuración
\end{itemize}

Cuando la depuración está habilitada, el firmware imprimirá los datos de control recibidos desde la aplicación Android.

Vuelve a cargar el firmware después de cualquier cambio.

\subsection{Paso 2: Abrir el Monitor Serie}

En el IDE de Arduino abre:

\begin{verbatim}
    Tools -> Serial Monitor
\end{verbatim}

Configura la velocidad a:

\begin{verbatim}
    115200
\end{verbatim}

Reinicia el ESP32 o presiona el botón de reset.

El Monitor Serie empezará a mostrar mensajes de depuración.

\subsection{Paso 3: Mover Controles en la App}

Después de conectar la aplicación Android al ESP32, mueve los controles en pantalla:

\begin{itemize}
    \item Mueve los joysticks
    \item Gira las perillas
    \item Activa o desactiva los interruptores
    \item Presiona botones
\end{itemize}

El ESP32 imprimirá los datos recibidos en el Monitor Serie.

\section{Tipos de Salida de Depuración}

El sistema de depuración puede imprimir diferentes categorías de información según la configuración en \texttt{debug\_config.h}.

\subsection{Depuración de Joystick (Sticks)}

Si la siguiente opción está habilitada:

\begin{verbatim}
    #define DBG_STICKS 1
\end{verbatim}

El ESP32 imprimirá las posiciones de los joysticks cuando cambien.

Ejemplo:

\begin{verbatim}
    L Stick X: 512    L Stick Y: 510
    L Stick X: 650    L Stick Y: 480

    R Stick X: 300    R Stick Y: 700
\end{verbatim}

Estos valores representan la posición del joystick izquierdo y derecho en la aplicación Android.

Rango típico:

\begin{itemize}
    \item Mínimo: alrededor de 0
    \item Centro: alrededor de 512
    \item Máximo: alrededor de 1023
\end{itemize}

Si los números cambian cuando mueves el joystick, la comunicación está funcionando correctamente.

\subsection{Depuración de Perillas (Knobs)}

Si la siguiente opción está habilitada:

\begin{verbatim}
    #define DBG_KNOBS 1
\end{verbatim}

El ESP32 imprimirá los valores de las perillas.

Ejemplo:

\begin{verbatim}
    Left Knob: 340
    Left Knob: 512
    Left Knob: 890

    Right Knob: 120
    Right Knob: 780
\end{verbatim}

Estos valores corresponden a la rotación de las perillas en la aplicación Android.

Las perillas pueden usarse para controlar parámetros como:

\begin{itemize}
    \item Velocidad del motor
    \item Ángulo del servo
    \item Intensidad de luz
    \item Inclinación/paneo de cámara
\end{itemize}

\subsection{Depuración de Interruptores (Switches)}

Si la siguiente opción está habilitada:

\begin{verbatim}
    #define DBG_SWITCHES 1
\end{verbatim}

El ESP32 mostrará el estado de los interruptores.

Ejemplo:

\begin{verbatim}
    Switch Byte: 0x05
    S1 = ON
    S2 = OFF
    S3 = ON
    S4 = OFF
    S5 = OFF
    S6 = OFF
\end{verbatim}

Cada interruptor corresponde a un botón tipo toggle en la aplicación Android.

Estos interruptores se usan comúnmente para:

\begin{itemize}
    \item Encender y apagar luces
    \item Activar relés
    \item Cambiar modos de conducción
    \item Habilitar funciones autónomas
\end{itemize}

\subsection{Depuración de Eventos}

Si la siguiente opción está habilitada:

\begin{verbatim}
    #define DBG_EVENTS 1
\end{verbatim}

El ESP32 imprimirá un mensaje cuando llegue un paquete de evento.

Ejemplo:

\begin{verbatim}
    --- EVENT PACKET (PRESS) ---
    Event ID: 0x03
    Checksum: 0x7A
    --------------
\end{verbatim}

Los eventos se disparan cuando el usuario presiona un botón en la aplicación Android.

Usos típicos:

\begin{itemize}
    \item Activar una función especial
    \item Reiniciar un sistema
    \item Disparar un sonido o animación
    \item Lanzar una acción en el robot
\end{itemize}

\section{Depuración del Sistema de Telemetría}

El firmware incluye un sistema de telemetría que permite al ESP32 enviar
información en tiempo real de vuelta a la aplicación Android.

La telemetría puede incluir:

\begin{itemize}
    \item Valores numéricos
    \item Indicadores
    \item Gráficas en tiempo real
\end{itemize}

Para desarrollo y pruebas, el firmware proporciona un sistema de depuración de telemetría
que permite simular valores mediante el Monitor Serie del IDE de Arduino.

Este sistema se configura en el archivo:

\begin{center}
    \texttt{debug\_config.h}
\end{center}

Dentro de este archivo puedes seleccionar el modo de depuración de telemetría:

\begin{verbatim}
    #define DBG_NONE        0
    #define DBG_PANEL       1
    #define DBG_INDICATOR   2
    #define DBG_PLOT        3

    #define TELEMETRY_DEBUG_MODE DBG_INDICATOR
\end{verbatim}

Solo debe estar activo un modo de depuración a la vez.

\subsection{Operación Normal}

\begin{verbatim}
    #define TELEMETRY_DEBUG_MODE DBG_NONE
\end{verbatim}

En este modo el ESP32 obtiene valores de telemetría desde sensores reales de hardware.

Ejemplos:

\begin{itemize}
    \item sensores analógicos
    \item medición de voltaje de batería
    \item entradas digitales
\end{itemize}

Este modo debe usarse cuando el robot o proyecto esté completamente armado.

\subsection{Modo Debug de Panel Numérico}

\begin{verbatim}
    #define TELEMETRY_DEBUG_MODE DBG_PANEL
\end{verbatim}

Este modo permite probar los paneles numéricos de telemetría sin conectar sensores.

Cuando este modo está habilitado, los valores ingresados en el Monitor Serie se interpretan y se envían
a la aplicación Android como telemetría de panel numérico.

\subsubsection{Formato de entrada en el Monitor Serie}

Abre el Monitor Serie del IDE de Arduino y envía valores con el siguiente formato:

\begin{verbatim}
    LEFT_VALUE,RIGHT_VALUE
\end{verbatim}

Ambos valores pueden incluir parte decimal.

Ejemplos:

\begin{verbatim}
    12.3,45.8
    120.0,350.5
    5,10
\end{verbatim}

\subsubsection{Conversión de valores}

Los valores ingresados en el Monitor Serie se convierten internamente al formato de telemetría usado por el firmware.

Cada valor se multiplica por 10 y se convierte a entero antes de transmitirse a la aplicación Android.

Ejemplos:

\begin{itemize}
    \item 12.3 se convierte en 123
    \item 45.8 se convierte en 458
    \item 120.0 se convierte en 1200
\end{itemize}

Del lado Android, el valor entero recibido se divide por 10 para mostrar el valor decimal correcto en el panel numérico.

\subsubsection{Rango válido}

El rango permitido para cada valor numérico es:

\begin{itemize}
    \item 0.0 a 6553.5
\end{itemize}

Este rango corresponde al valor máximo transmisible usando un entero sin signo de 16 bits con una cifra decimal de precisión.

\subsubsection{Ejemplo}

Si se ingresa el siguiente valor en el Monitor Serie:

\begin{verbatim}
    123.4,56.7
\end{verbatim}

El firmware transmitirá:

\begin{verbatim}
    1234,567
\end{verbatim}

La aplicación Android mostrará:

\begin{verbatim}
    123.4   56.7
\end{verbatim}

en los paneles numéricos izquierdo y derecho respectivamente.

Este modo es útil para verificar la transmisión de telemetría y la interfaz Android antes de conectar sensores reales.

\subsection{Modo Debug de Indicadores}

\begin{verbatim}
    #define TELEMETRY_DEBUG_MODE DBG_INDICATOR
\end{verbatim}

Este modo permite probar widgets de indicadores.

Envía valores por el Monitor Serie usando este formato:

\begin{verbatim}
    ANALOG_VALUE,BATTERY_LEVEL,DIGITAL_MASK
\end{verbatim}

Ejemplo:

\begin{verbatim}
    45,80,5
\end{verbatim}

Rangos:

\begin{itemize}
    \item Valor analógico: 0 a 100
    \item Nivel de batería: 0 a 100
    \item Máscara digital: 0 a 255
\end{itemize}

La máscara digital representa el estado de indicadores digitales usando flags por bits.

Ejemplos:

\begin{verbatim}
    1
    2
    4
    8
\end{verbatim}

Cada valor activa un indicador diferente.

Se pueden activar múltiples indicadores combinando valores.

Ejemplo:

\begin{verbatim}
    3
\end{verbatim}

Representación binaria:

\begin{verbatim}
    00000011
\end{verbatim}

Esto activa los dos primeros indicadores.

\subsection{Modo Debug de Gráficas (Plot)}

\begin{verbatim}
    #define TELEMETRY_DEBUG_MODE DBG_PLOT
\end{verbatim}

Este modo simula tres señales de telemetría que se muestran como gráficas en tiempo real en la aplicación Android.

Envía valores usando el Monitor Serie:

\begin{verbatim}
    VALUE1,VALUE2,VALUE3
\end{verbatim}

Ejemplo:

\begin{verbatim}
    10,40,80
\end{verbatim}

Rango válido:

\begin{itemize}
    \item 0 a 255
\end{itemize}

Estos valores aparecerán en las gráficas de telemetría de la aplicación.

\subsection{Cómo Funciona el Sistema de Depuración de Telemetría}

Los valores de depuración son leídos por el módulo:

\begin{center}
    \texttt{telemetry\_source.cpp}
\end{center}

Este módulo reemplaza valores de sensores reales por valores recibidos vía interfaz serie cuando hay un modo debug activo.

Los paquetes de telemetría se crean y transmiten mediante:

\begin{center}
    \texttt{telemetry.cpp}
\end{center}

Este módulo arma paquetes de telemetría y los envía a la aplicación Android usando Bluetooth.

Esta arquitectura permite probar la interfaz de telemetría sin conectar sensores físicos.

\section{Por Qué el Modo Debug es Importante}

El modo debug permite verificar rápidamente que:

\begin{itemize}
    \item La aplicación Android está conectada
    \item Se están recibiendo datos por Bluetooth
    \item Los valores de control cambian correctamente
\end{itemize}

Si los valores aparecen en el Monitor Serie pero el hardware no responde, el problema probablemente está en el cableado de hardware o en la configuración del firmware, y no en la conexión Bluetooth.

\section{Personalizar Etiquetas de Telemetría en la Aplicación Android}

La aplicación Android muestra automáticamente etiquetas para valores de telemetría
como paneles numéricos, indicadores y gráficas.

Estas etiquetas se definen en el firmware del ESP32 y se envían a la aplicación
cuando se establece la conexión Bluetooth.

Esto permite personalizar la interfaz sin modificar la aplicación Android.

\subsection{Dónde se Definen las Etiquetas}

Las etiquetas de telemetría se definen dentro del archivo:

\begin{center}
    \texttt{telemetry.cpp}
\end{center}

Dentro de la función encargada de enviar la configuración de telemetría, encontrarás estas variables:

\begin{verbatim}
    const char* plotNames[] = { "Volts", "Amps", "RPMs" };
    const char* panelNames[] = { "Left", "Right" };
    const char* indicatorNames[] = { "Throttle", "Battery" };
\end{verbatim}

Estos arreglos definen las etiquetas que aparecen en la pantalla RC de la aplicación Android.

\subsection{Etiquetas de Gráficas (Plot)}

El arreglo \texttt{plotNames} define los nombres de las gráficas en tiempo real.

Ejemplo:

\begin{verbatim}
    const char* plotNames[] = { "Volts", "Amps", "RPMs" };
\end{verbatim}

En la aplicación Android esto creará tres gráficas con las etiquetas:

\begin{itemize}
    \item Volts
    \item Amps
    \item RPMs
\end{itemize}

Puedes renombrarlas para que coincidan con tus sensores.

Ejemplo para un robot tipo auto:

\begin{verbatim}
    const char* plotNames[] = { "Speed", "Motor Temp", "Battery Current" };
\end{verbatim}

\subsection{Etiquetas de Panel Numérico}

El arreglo \texttt{panelNames} define las etiquetas de los paneles numéricos grandes.

Ejemplo:

\begin{verbatim}
    const char* panelNames[] = { "Left", "Right" };
\end{verbatim}

Usos típicos:

\begin{itemize}
    \item Velocidad de motor
    \item Distancia
    \item RPM
    \item Temperatura
\end{itemize}

Ejemplo de modificación:

\begin{verbatim}
    const char* panelNames[] = { "Motor RPM", "Battery Voltage" };
\end{verbatim}

Nota: se recomienda usar un máximo de 12 caracteres.

\subsection{Etiquetas de Indicadores}

El arreglo \texttt{indicatorNames} define las etiquetas de los widgets de indicadores.

Ejemplo:

\begin{verbatim}
    const char* indicatorNames[] = { "Throttle", "Battery" };
\end{verbatim}

Posibles usos:

\begin{itemize}
    \item Posición del acelerador
    \item Nivel de batería
    \item Temperatura
    \item Intensidad de señal
\end{itemize}

Ejemplo de modificación:

\begin{verbatim}
    const char* indicatorNames[] = { "Engine Load", "Battery Level" };
\end{verbatim}

\subsection{Cuándo se Envían las Etiquetas}

Las etiquetas de telemetría se transmiten solo una vez después de que la aplicación Android se conecta al ESP32.

El firmware envía un paquete de configuración de telemetría que incluye:

\begin{itemize}
    \item etiquetas de gráficas
    \item etiquetas de paneles numéricos
    \item etiquetas de indicadores
\end{itemize}

Después de recibir este paquete, la aplicación Android actualiza automáticamente la interfaz de la pantalla RC.

\subsection{Notas Importantes}

\begin{itemize}
    \item El número de etiquetas debe coincidir con el número de variables de telemetría.
    \item Después de modificar las etiquetas, debes volver a cargar el firmware.
    \item La aplicación Android actualizará las etiquetas automáticamente tras reconectar.
\end{itemize}